\documentclass[11pt]{article}
\usepackage[margin=1in]{geometry}
\usepackage{amsmath,amssymb,amsthm}
\usepackage{graphicx}
\usepackage[hidelinks]{hyperref}
\setlength{\emergencystretch}{2em}

\newtheorem{theorem}{Theorem}
\newtheorem{lemma}{Lemma}
\newtheorem{definition}{Definition}
\newtheorem{proposition}{Proposition}

\newcommand{\C}{\mathbb C}
\newcommand{\Kcal}{\mathcal K}
\newcommand{\Prim}{\operatorname{Prim}}
\newcommand{\Ann}{\operatorname{Ann}}
\newcommand{\id}{\operatorname{id}}

\title{Finite-Dimensional Quotients of Separable-Fibre Continuous-Trace Algebras Have BCP}
\author{Run \texttt{fa\_banach\_001}, lane 19}
\date{June 17, 2026}

\begin{document}
\maketitle

\section*{Source Problem}

The source paper is Jinghao Huang, Karimbergen Kudaybergenov, and Rui Liu,
\emph{The ball-covering property of von Neumann algebras and noncommutative
symmetric spaces}, arXiv:2311.14257.  After proving that a commutative unital
\(C^*\)-algebra \(C(K)\) has BCP/UBCP exactly when \(K\) has a countable
\(\pi\)-basis, the authors ask for a noncommutative version:

\begin{center}
\includegraphics[width=0.92\textwidth]{figures/open_problem_crop.png}
\end{center}

\section*{Definitions}

\begin{definition}
A \(C^*\)-algebra \(A\) has the ball-covering property, or BCP, if the unit
sphere of \(A\) is covered by countably many open balls none of which contains
the origin.  It has the uniform ball-covering property, or UBCP, if the balls
can be chosen with a common positive gap from the origin: for centers
\(c_n\) and radii \(r_n\),
\[
        \inf_n(\|c_n\|-r_n)>0.
\]
\end{definition}

\begin{definition}
Let \(J\) be a closed ideal of a \(C^*\)-algebra \(A\).  We say that \(J\) is
\(\gamma\)-norming in \(A\), where \(0<\gamma\le1\), if
\[
        \sup\{\|\pi(a)\|:\pi\in\widehat J\}\ge \gamma\|a\|
        \qquad(a\in A).
\]
\end{definition}

\section*{Statement}

This packet removes the norming hypothesis from the previous finite-quotient
packet.  The result is still a partial answer to the original source problem,
because the ideal is assumed continuous trace and the quotient is finite
dimensional.

\begin{theorem}[finite-dimensional quotient continuous-trace criterion]
\label{thm:main}
Let
\[
        0\longrightarrow I\longrightarrow A
        \stackrel{q}{\longrightarrow} F\longrightarrow0
\]
be an exact sequence of \(C^*\)-algebras.  Assume that \(F\) is finite
dimensional and that \(I\) is continuous trace with elementary fibres
\(\Kcal(H_x)\), where each \(H_x\) is separable and nonzero.

If \(\Prim(I)\) has a countable \(\pi\)-basis, then \(A\) has UBCP.
Conversely, if \(A\) has BCP, then \(\Prim(I)\) has a countable \(\pi\)-basis.
Consequently, in this finite-dimensional quotient continuous-trace class,
\[
        A\text{ has BCP}
        \quad\Longleftrightarrow\quad
        A\text{ has UBCP}
        \quad\Longleftrightarrow\quad
        \Prim(I)\text{ has a countable }\pi\text{-basis}.
\]
\end{theorem}

\section*{Idea of the Proof}

The earlier matrix-block ballast theorem proves the result under the extra
assumption that \(I\) is uniformly norming in \(A\).  At first the missing case
seems to require a separate quotient-point cover: perhaps some elements are
almost invisible on \(\Prim(I)\) and are normed mainly by the finite quotient.

The key observation is that this is not a new covering problem.  The invisible
part is exactly the annihilator
\[
        J=\Ann_A(I)=\{a\in A:aI=Ia=0\}.
\]
The Busby map \(\tau:F\to Q(I)=M(I)/I\) has kernel \(q(J)\).  Thus, after
splitting off the finite-dimensional ideal \(J\) and replacing \(I\) by
\[
        I^\sharp=I\oplus J,
\]
the remaining quotient has injective Busby map.  Injective \(C^*\)-homomorphisms
are isometric, so the action on the ideal sees the quotient norm.  This makes
\(I^\sharp\) uniformly norming in \(A\), and the previous norming theorem
applies.

\section*{Imported Norming Theorem}

We use the following result from the packet
\texttt{2311.14257\_norming\_finite\_quotient\_bcp}.

\begin{proposition}[norming finite-dimensional quotient criterion]
\label{prop:norming}
Let
\[
        0\longrightarrow B\longrightarrow E\longrightarrow G\longrightarrow0
\]
be exact.  Assume that \(G\) is finite dimensional, \(B\) is continuous trace
with separable nonzero elementary fibres, \(\Prim(B)\) has a countable
\(\pi\)-basis, and \(B\) is \(\gamma\)-norming in \(E\) for some
\(\gamma>0\).  Then \(E\) has UBCP.
\end{proposition}

The proof of Proposition~\ref{prop:norming} is the matrix-block ballast
argument: choose a bounded linear section of the finite quotient, approximate
the compact ideal component locally, and insert a rank-one spike aligned with
the local quotient-lifted operator.  The proof does not require a
\(*\)-splitting of the quotient.

\section*{Proof}

We first isolate the exact obstruction to norming.

\begin{lemma}[Busby kernel equals the quotient of the annihilator]
\label{lem:busby-kernel}
Let
\[
        0\longrightarrow I\longrightarrow A
        \stackrel{q}{\longrightarrow} F\longrightarrow0
\]
be exact, and let \(\phi:A\to M(I)\) be the canonical action homomorphism.
Put
\[
        J=\Ann_A(I)=\{a\in A:aI=Ia=0\}.
\]
Let \(\tau:F\to Q(I)=M(I)/I\) be the Busby map.  Then
\[
        \ker\phi=J,
        \qquad
        \ker\tau=q(J).
\]
In particular, if \(F\) is finite dimensional, then \(J\) is finite
dimensional.
\end{lemma}

\begin{proof}
By definition, \(\phi(a)\) is the multiplier of \(I\) obtained by left and
right multiplication by \(a\).  Hence \(\phi(a)=0\) exactly when \(aI=Ia=0\),
so \(\ker\phi=J\).

If \(j\in J\), then \(\phi(j)=0\), and therefore
\[
        \tau(q(j))=\phi(j)+I=0.
\]
Thus \(q(J)\subseteq\ker\tau\).

Conversely, suppose \(\tau(q(a))=0\).  Then \(\phi(a)\) belongs to the copy of
\(I\) inside \(M(I)\).  Hence there is \(i\in I\) such that
\(\phi(a)=\phi(i)\).  Therefore \(a-i\in\ker\phi=J\), and
\[
        q(a)=q(a-i)\in q(J).
\]
So \(\ker\tau=q(J)\).

Finally \(J\cap I=\{0\}\): if \(x\in J\cap I\), then \(x x^*=0\) because
\(x^*\in I\) and \(xI=0\), hence \(x=0\).  Thus \(q|_J\) is injective.  Since
\(q(J)\subset F\), finite dimensionality of \(F\) implies finite dimensionality
of \(J\).
\end{proof}

\begin{lemma}[finite annihilator split]
\label{lem:split}
In the setting of Lemma~\ref{lem:busby-kernel}, assume \(F\) is finite
dimensional.  Then \(J=\Ann_A(I)\) is a finite-dimensional unital ideal.  If
\(p\in J\) is its identity, then \(p\) is central in \(A\), and
\[
        A=J\oplus A_0,
        \qquad
        A_0=\{a\in A:pa=0\}.
\]
Moreover \(I\subset A_0\), and \(A_0/I\cong F/q(J)\).
\end{lemma}

\begin{proof}
A finite-dimensional \(C^*\)-algebra is unital; let \(p\) be the unit of
\(J\).  For any \(a\in A\), both \(pa\) and \(ap\) belong to the ideal \(J\).
Since \(p\) is the identity on \(J\),
\[
        pa=(pa)p=p(ap)=ap.
\]
Thus \(p\) is central in \(A\).

Every \(a\in A\) decomposes as
\[
        a=pa+(a-pa),
\]
with \(pa\in J\) and \(p(a-pa)=0\).  The intersection
\(J\cap A_0\) is zero because \(p\) acts as the identity on \(J\).  Hence
\(A=J\oplus A_0\) as a \(C^*\)-direct sum.  Since \(J\) annihilates \(I\), we
have \(pI=0\), so \(I\subset A_0\).

The map \(A_0\to F/q(J)\), \(a_0\mapsto q(a_0)+q(J)\), is onto because
\[
        q(a)=q(a-pa)+q(pa),
        \qquad
        q(pa)\in q(J).
\]
Its kernel is \(I\): if \(a_0\in A_0\) and \(q(a_0)\in q(J)\), then
\(a_0-j\in I\) for some \(j\in J\).  Applying \(p\) gives
\[
        0-j=p(a_0-j)\in I,
\]
so \(j\in I\cap J=\{0\}\), and \(a_0\in I\).
\end{proof}

\begin{lemma}[the enlarged ideal is norming]
\label{lem:enlarged-norming}
Assume the hypotheses of Theorem~\ref{thm:main}.  Let
\[
        J=\Ann_A(I),
        \qquad
        I^\sharp=I\oplus J.
\]
Then \(I^\sharp\) is a continuous-trace ideal of \(A\) with separable nonzero
elementary fibres.  If \(\Prim(I)\) has a countable \(\pi\)-basis, then so
does \(\Prim(I^\sharp)\).  Moreover \(I^\sharp\) is \(1\)-norming in \(A\).
\end{lemma}

\begin{proof}
By Lemma~\ref{lem:split}, \(A=J\oplus A_0\), \(I\subset A_0\), and
\[
        0\longrightarrow I\longrightarrow A_0
        \longrightarrow F/q(J)\longrightarrow0
\]
is exact.  The Busby map of this reduced extension is the quotient of
\(\tau\) by its kernel \(q(J)\), hence is injective by
Lemma~\ref{lem:busby-kernel}.  Since it is an injective \(C^*\)-homomorphism,
it is isometric.

Let \(a_0\in A_0\), and put
\[
        y=a_0+I\in A_0/I\cong F/q(J).
\]
Let \(\phi_0:A_0\to M(I)\) be the canonical action.  The norm of
\(\phi_0(a_0)\) is the supremum of the norms of \(a_0\) in the primitive
representations over \(\Prim(I)\).  Write this number as
\[
        \alpha(a_0)
        =
        \sup\{\|\pi(a_0)\|:\pi\in\widehat I\}.
\]
The Busby value satisfies
\[
        \|y\|
        =
        \|\tau_0(y)\|
        =
        \|\phi_0(a_0)+I\|_{Q(I)}
        \le
        \|\phi_0(a_0)\|
        =
        \alpha(a_0).
\]
The norm in \(A_0\) is the supremum over primitive representations that either
restrict nontrivially to \(I\), or factor through \(A_0/I\).  Therefore
\[
        \|a_0\|=\max\{\alpha(a_0),\|y\|\}=\alpha(a_0).
\]
Thus \(I\) is already \(1\)-norming in \(A_0\).

Now take an arbitrary \(a\in A\), and write \(a=j+a_0\) with \(j\in J\) and
\(a_0\in A_0\).  Because \(A=J\oplus A_0\) is a \(C^*\)-direct sum,
\[
        \|a\|=\max\{\|j\|,\|a_0\|\}.
\]
The primitive representations of \(I^\sharp=J\oplus I\) see the two summands,
so
\[
\begin{aligned}
        \sup_{\rho\in\widehat{I^\sharp}}\|\rho(a)\|
        &=
        \max\left\{
              \|j\|,
              \sup_{\pi\in\widehat I}\|\pi(a_0)\|
             \right\}  \\
        &=
        \max\{\|j\|,\|a_0\|\}
        =
        \|a\|.
\end{aligned}
\]
Hence \(I^\sharp\) is \(1\)-norming in \(A\).

Finally, \(I^\sharp=I\oplus J\) is a direct sum of the given continuous-trace
algebra and a finite-dimensional \(C^*\)-algebra, hence is continuous trace.
Its elementary fibres are those of \(I\) together with finitely many matrix
algebras, all over separable nonzero Hilbert spaces.  Also
\[
        \Prim(I^\sharp)=\Prim(I)\sqcup\Prim(J)
\]
as a topological disjoint union, and \(\Prim(J)\) is finite discrete.  Thus a
countable \(\pi\)-basis for \(\Prim(I)\), together with the finitely many
singleton points of \(\Prim(J)\), gives a countable \(\pi\)-basis for
\(\Prim(I^\sharp)\).
\end{proof}

\begin{proof}[Proof of Theorem~\ref{thm:main}]
Assume first that \(\Prim(I)\) has a countable \(\pi\)-basis.  By
Lemma~\ref{lem:enlarged-norming}, the enlarged ideal
\[
        I^\sharp=I\oplus\Ann_A(I)
\]
is continuous trace with separable nonzero elementary fibres, has a countable
\(\pi\)-basis on its primitive spectrum, and is \(1\)-norming in \(A\).  The
quotient
\[
        A/I^\sharp\cong F/q(\Ann_A(I))
\]
is finite dimensional.  Proposition~\ref{prop:norming}, applied to
\[
        0\longrightarrow I^\sharp\longrightarrow A
        \longrightarrow A/I^\sharp\longrightarrow0,
\]
therefore gives UBCP for \(A\).

Conversely, if \(A\) has BCP, the primitive \(\pi\)-base obstruction for
\(C^*\)-algebras from the earlier lane 19 obstruction packet gives a
countable \(\pi\)-basis for \(\Prim(A)\).  Since
\(\Prim(I)\) is an open subspace of \(\Prim(A)\), it inherits a countable
\(\pi\)-basis.  This proves the converse implication.  The equivalence follows
because UBCP implies BCP.
\end{proof}

\section*{Scope and Limitations}

This is a genuine removal of the norming hypothesis for finite-dimensional
quotients over the positive continuous-trace/separable-fibre class.  It is not
claimed as a full solution to the source paper's broad noncommutative problem:
the proof still uses continuous trace structure and a finite-dimensional
quotient.

The next natural targets are:
\begin{enumerate}
    \item finite composition series whose subquotients are of the above form;
    \item Fell or type-I algebras with separable irreducible representations;
    \item a general obstruction theorem showing when non-type-I primitive
          fibres force failure of BCP.
\end{enumerate}

\section*{Novelty and Search Bounds}

This packet was produced as a strengthening of the earlier lane 19 norming
finite-quotient packet.  During the surrounding run, the local indexes were
searched for arXiv:2311.14257, finite-dimensional
quotients, finite-dimensional extensions, unitization, matrix-block ballast,
continuous trace, Busby maps, annihilator ideals, BCP, UBCP, and primitive
\(\pi\)-bases.  Earlier web searches on June 17, 2026 for BCP finite
codimensional permanence, finite-dimensional \(C^*\)-quotient splitting, and
Fell/type-I local continuous-trace structure found the source paper,
vector-valued \(C_0(K,X)\) BCP stability arXiv:2111.04921,
Calkin/operator-space BCP papers arXiv:2412.07137 and arXiv:2507.02261, and
Fell-algebra structure arXiv:1207.2540, but no exact finite-dimensional
quotient continuous-trace BCP theorem.  No claim of complete novelty is made
without human literature review.

\section*{Human Review Focus}

The proof reduces the new theorem to one standard extension-theoretic
calculation.  The most important points to check are:
\begin{enumerate}
    \item the equality \(\ker\tau=q(\Ann_A(I))\);
    \item centrality of the unit of the finite-dimensional annihilator ideal;
    \item the assertion that the reduced Busby map is injective and therefore
          isometric;
    \item the norm formula
          \(\|a_0\|=\max\{\alpha(a_0),\|a_0+I\|\}\) for the reduced extension.
\end{enumerate}
If these pass review, the finite-dimensional quotient obstruction from the
previous packet is removed.

\begin{thebibliography}{9}

\bibitem{HKL}
Jinghao Huang, Karimbergen Kudaybergenov, and Rui Liu,
\emph{The ball-covering property of von Neumann algebras and noncommutative
symmetric spaces},
arXiv:2311.14257, 2023.

\bibitem{lane19norming}
Run \texttt{fa\_banach\_001}, lane 19,
\emph{A Norming Finite-Quotient Criterion for \(C^*\)-Algebras with UBCP},
solution packet
\texttt{2311.14257\_norming\_finite\_quotient\_bcp}, 2026.

\bibitem{lane19unit}
Run \texttt{fa\_banach\_001}, lane 19,
\emph{Unitizations of Separable-Fibre Continuous-Trace Algebras Have BCP},
solution packet
\texttt{2311.14257\_unitization\_continuous\_trace\_bcp}, 2026.

\bibitem{lane19pi}
Run \texttt{fa\_banach\_001}, lane 19,
\emph{Primitive \(\pi\)-Base Obstruction for \(C^*\)-Algebras with BCP},
solution packet
\texttt{2311.14257\_primitive\_pi\_base\_obstruction}, 2026.

\bibitem{Dixmier}
J. Dixmier,
\emph{\(C^*\)-Algebras},
North-Holland, 1977.

\bibitem{RW}
I. Raeburn and D. P. Williams,
\emph{Morita Equivalence and Continuous-Trace \(C^*\)-Algebras},
American Mathematical Society, 1998.

\end{thebibliography}

\end{document}
